\documentclass[11pt]{article}
\usepackage[a4paper,margin=24mm]{geometry}
\usepackage{amsmath,amssymb,amsthm,booktabs}
\newtheorem{lemma}{Lemma}\newtheorem{theorem}{Theorem}
\title{A Baer-subplane construction of Geometry 3.3\\and a correction to its printed incidence rule}
\author{L100 finite construction component}\date{}
\begin{document}\maketitle
We use the fixed 22-page Halle preprint of Gottschalk and Leemans
\cite{GL}, physical page 19. Its later journal edition has not been
collated. Put $P=\operatorname{PG}(2,4)$ and $G=\operatorname{PSL}_3(4)$.
The coordinate plane, its exact identification with the retained
21-point action, and the complete coset representative numbered 10 are
acquired from the accepted projective-construction, palette and source-join
components. The latter representative is Geometry 3.3.

\begin{theorem}
For each of the three $G$-orbits of Baer subplanes, the three incidence
rules printed for Geometry 3.3 admit a two-object flag which does not
extend to a chamber. Retain the first two rules and restrict the third
to pairs having a common incident Baer subplane in the selected orbit.
Each of the three resulting incidence systems is isomorphic, with its
specified $G$-action transported by an explicit automorphism of $G$, to
the acquired Geometry 3.3. The printed and corrected line-pair to
point-pair degrees are respectively $18$ and $6$.
\end{theorem}
This corrects an ancillary construction recipe. It does not refute the
paper's principal classification, which already contains the corrected
geometry. Other constructions and the complete source disposition are
outside this component.

\section{Complete subplane input}
\begin{lemma}
A seven-point subset $B$ of $P$ is an embedded plane of order two if
and only if every ambient line meets it in one or three points.
There are exactly $360$ such subsets.
\end{lemma}
\begin{proof}
Under the intersection condition every pair of points lies on a unique
three-point intersection. These give $\binom72/\binom32=7$ lines,
three through each point. Each such line meets the six others, twice
at each of its three points, since two lines cannot share two points.
Thus the induced incidence system is an embedded Fano plane.
Conversely, a Fano subplane has seven ambient lines meeting it in three
points. Any other ambient line meets it in at most one point. It cannot
avoid the subplane: its intersections with the seven subplane lines
would be distinct, although it has only five points.

An ordered projective frame is a sequence of four points with no three
collinear. There are $21\cdot20\cdot16\cdot9=60480$ ambient frames.
Every frame lies in a unique Fano subplane. Indeed a projectivity sends
it to $[1,0,0],[0,1,0],[0,0,1],[1,1,1]$. The three remaining Fano points
are forced by intersections of the joins of these four points; they
are $[1,1,0],[1,0,1],[0,1,1]$. These seven points do form a subplane,
so existence and uniqueness both hold. Each Fano plane has
$7\cdot6\cdot4\cdot1=168$ ordered frames. Counting frames with their
unique subplane gives $60480/168=360$.
\end{proof}
The preserved first finite pass examined all $\binom{21}{7}=116280$
subsets and saved the $360$ qualifying sets before failing its later,
uncorrected incidence join. The current checker uses these saved sets:
it verifies distinctness and the lemma's criterion, then closes each
listed orbit under the actual acquired generators. The result is three
orbits of size $120$. Completeness follows from the lemma; it is not
inferred from a printed orbit count. No subset search is repeated.

\section{The printed rule and its repair}
For one orbit $\mathcal B$, the types are $\mathcal B$, unordered pairs
of distinct lines, and unordered pairs of distinct points. Write $pq$
for the line joining $p$ and $q$. The printed rules are
\begin{align}
B*\{\ell,m\}&\iff B\cap(\ell\cup m)=\ell\cap m,\tag{A}\\
B*\{p,q\}&\iff B\cap\{p,q\}=\varnothing\text{ and }|B\cap pq|=3,\tag{B}\\
\{\ell,m\}*\{p,q\}&\iff (\ell\cup m)\cap\{p,q\}=\varnothing
\text{ and }\ell\cap m\cap pq\ne\varnothing.\tag{C}
\end{align}
These are literal transcriptions of the native source, not of its OCR.
Let $o=\ell\cap m$. In (C), $pq$ must be one of the three other lines
through $o$. On each there are four points outside $\ell\cup m$.
Consequently (C) has degree $3\binom42=18$ at every line pair.

Here is an explicit finite obstruction to flag extension. Label the
retained points $0,\ldots,20$ as in the acquired coordinate-to-permutation
map, and encode a point set $X$ by $\sum_{x\in X}2^x$. Fix
$\ell=\{0,1,2,3,4\}$ and $m=\{0,5,6,7,8\}$. In the three orbits
ordered by their least encoded subplane, the following table gives a
printed incident point pair and the complete two sets of Baer neighbours.
\begin{center}\small
\begin{tabular}{ccll}\toprule
Orbit&Point pair&Neighbours of $\{\ell,m\}$&Neighbours of the point pair\\\midrule
0&$\{9,11\}$&710657, 730625,&46093, 87059,\\
&&1361409, 1390593&398529, 1578273\\
1&$\{9,10\}$&816129, 887297,&47297, 88353,\\
&&1206785, 1283073&399373, 1579027\\
2&$\{9,10\}$&438785, 478209,&30741, 104459,\\
&&1616897, 1659393&792737, 1186113\\\bottomrule
\end{tabular}
\end{center}
The joining line in each case is $\{0,9,10,11,12\}$, so (C) holds.
The two neighbour sets in each row are disjoint. The checker obtains
each list by applying (A) and (B) to every member of that orbit, using
an independent set formulation. Thus each displayed pair is a flag
with empty remaining-type residue. This proves the stated defect of
the printed recipe.

The repaired relation is
\[
\{\ell,m\}*'\{p,q\}\iff\text{(C) holds and there exists }B\in\mathcal B
\text{ satisfying (A) and (B).}
\]
Rules (A) and (B) are unchanged. The selected orbit is essential.
For additional exact context, the $3780$ printed pairs split into three
sets of $1260$: in each pair precisely one orbit supplies common Baer
neighbours, and it supplies two. The checker verifies this finite
partition directly. In each repaired system the degree is six.

\section{Literal positive identifications}
The acquired catalogue representative has maximal parabolics $H_i$
and objects the right cosets $H_i g$. For each of the three repaired
systems the certificate supplies a type permutation, base objects
$x_i$ and a point permutation $h$. It checks that $h$ preserves the
actual lines and conjugates the full, explicit $20160$-element group
onto itself. With $\alpha(g)=h^{-1}gh$, the object map is
\[
             H_i g\longmapsto x_i^{\alpha(g)}.
\]
The literal parabolic generators fix $x_i$. Moreover the verifier
checks all $60480$ coset-member assignments in each map, so it verifies
well-definedness independently of the choice of coset representative.
It then checks each typewise bijection, all $189000$ directed
cross-type object pairs, and all $1080$ generator-object action pairs.
Line pairs are pairs of five-point sets acted on by actual point
permutations; no point-label/line-label identification is assumed.

All three maps pass. Therefore incidence, nonincidence and the specified
action agree exactly. The already proved geometry, connectedness,
residual properties, flag transitivity and RWPRI properties transport
through these maps. In particular the three acquired local
two-transitivity flags are all true. Equivalently, composing one map
with the inverse of another gives a full incidence equivalence between
any two orbit choices; this is not inferred from their equal sizes.

The certificate's finite verification, its exact acquired input
interfaces and the short lemmas above constitute the proof. The
acquired group order and full coset catalogue retain their previously
stated computational and classification premises. The positive maps
do not require a new group, subgroup, automorphism or graph-isomorphism
search. The first failed join, the prelaunch ordering correction and
the changed successful join are preserved separately. This component
makes no novelty, measurement, open-question or whole-source claim.

\begin{thebibliography}{1}
\bibitem{GL} H. Gottschalk and D. Leemans,
\emph{Geometries for the group $\operatorname{PSL}(3,4)$},
retained Halle preprint 99--30 (1999), physical page 19.
The later journal edition is not collated here.
\end{thebibliography}
\end{document}
